\documentclass[11pt, a4paper]{article}
\usepackage{geometry}
\usepackage[parfill]{parskip}
\usepackage{graphicx}
\usepackage{amsmath}
\usepackage{enumerate}

\usepackage{charter}

\title{Huiswerk 5 \\ IB0402 Logica, verzamelingen en relaties}
\author{Harman Singh}
\date{Month Day, Year}

\begin{document}
\maketitle

\pagebreak
\section*{Opgave 14}
\subsection*{oefening a}

Om te bewijzen dat R een equivalentierelatie is, moet ik aantonen dat R reflexief, symmetrisch en transitief is.

\textbf{Reflexief}

Een relatie is reflexief als $\forall a \in A$ geldt dat $aRa$. Dit betekent dat elk reëel getal in een relatie staat met zichzelf.

Met de notatie die we hier gebruiken zal dat altijd het geval zijn, omdat de bewerking die je uitvoert enkel 1 antwoord kan hebben: $\lfloor 4.01 \rfloor$ is altijd 4, en kan niks anders zijn. Bij $aRa$ voer je twee keer dezelfde bewerking uit, dus voor elk reëel getal $a$ geldt dat $aRa$.


\textbf{Symmetrisch}

Een relatie is symmetrisch als $\forall a, b \in A$ geldt dat als $aRb$, dan ook $bRa$. Dit betekent dat als een reëel getal in een relatie staat met een ander reëel getal, dat dat andere reëel getal ook in een relatie staat met het eerste reeel getal.

Om dit te bewijzen zal ik 2 situaties bekijken
\begin{itemize}
	\item twee verschillende negatieve reële getallen, en
	\item twee verschillende positieve reële getallen
\end{itemize}

In de opgave staan reeds twee positieve reële getallen die beide hetzelfde resultaat opleveren bij de bewerking: 4.01 en 4.77.
Hieruit volgt dat ook $\lfloor 4.77 \rfloor R \lfloor 4.01 \rfloor$ geldt en ook $\lfloor 4.01 \rfloor R \lfloor 4.77 \rfloor$, omdat beide getallen dezelfde uitkomst opleveren bij de bewerking, namelijk $4$. 

Voor de negatieve reële getallen neem ik -2.1 en -2.9. Als je voor -2.1 het grootste geheel getal zoekt dat kleiner is dan -2.1, dan kom je aan $-3$. En als je dat voor -2.9 doet, dan kom je ook aan $-3$. Dus ook hier geldt dat $\lfloor -2.1 \rfloor R \lfloor-2.9 \rfloor$ en $\lfloor -2.9 \rfloor R \lfloor -2.1 \rfloor$.

\textbf{Transitief}

Hier neem ik drie verschillende positieve reële getallen: 5.23, 5.01 en 5.71, en reken ik de bewerking uit.

$\lfloor 5.23 \rfloor$ = 5; $\lfloor 5.01 \rfloor$ = 5; $\lfloor 5.71 \rfloor$ = 5

Hier geldt dus dat $\lfloor 5.23 \rfloor R \lfloor 5.01 \rfloor$ en $\lfloor 5.01 \rfloor R \lfloor 5.71 \rfloor$, en dus ook dat $\lfloor 5.23 \rfloor R \lfloor 5.71 \rfloor$.

\subsection*{oefening b}
\ldots


\section*{Opgave 15}
\subsection*{oefening a}
\ldots


\subsection*{oefening b}
\ldots


\subsection*{oefening c}
\ldots


\subsection*{oefening d}
\ldots


\subsection*{oefening e}
\ldots







\ldots
\end{document}